% 基线JPEG编码流程；只画数据流，不表示各步的计算量。宽约150mm，跨栏使用。
\begin{tikzpicture}[font=\figurefont\small,
  stage/.style={draw=FigureBlue,fill=FigureBlueFill,rounded corners=1.8pt,
    minimum width=21mm,minimum height=9mm,align=center,inner sep=1.5mm},
  lossy/.style={stage,draw=FigureRed,fill=FigureRedFill},
  code/.style={stage,draw=FigureYellow,fill=FigureYellowFill},
  flow/.style={-{Stealth[length=1.8mm]},draw=BookMuted,line width=0.8pt},
  note/.style={text=BookMuted,font=\figurefont\footnotesize}]
 \node[stage] (rgb) at (0,0) {RGB图像};
 \node[lossy] (ycc) at (3.0,0) {YCbCr变换\\色度下采样};
 \node[stage] (blk) at (6.0,0) {$8\times8$分块\\减去128};
 \node[stage] (dct) at (9.0,0) {二维DCT};
 \node[lossy] (q) at (12.0,0) {除以量化表\\取整};
 \node[stage] (zz) at (12.0,-1.75) {之字形扫描};
 \node[code] (dc) at (9.0,-1.1) {DC：与前块\\作差分};
 \node[code] (ac) at (9.0,-2.4) {AC：零游程\\与幅值配对};
 \node[code] (hf) at (6.0,-1.75) {Huffman\\变长编码};
 \node[stage] (bs) at (3.0,-1.75) {码流\\含表与头信息};
 \draw[flow] (rgb)--(ycc);
 \draw[flow] (ycc)--(blk);
 \draw[flow] (blk)--(dct);
 \draw[flow] (dct)--(q);
 \draw[flow] (q)--(zz);
 \draw[flow] (zz)--(dc);
 \draw[flow] (zz)--(ac);
 \draw[flow] (dc)--(hf);
 \draw[flow] (ac)--(hf);
 \draw[flow] (hf)--(bs);
 \node[note] at (7.5,.85) {红框为有损步骤，其余步骤可逆或无损};
\end{tikzpicture}
